% Part: methods
% Chapter: induction
% Section: inductive-definitions

\documentclass[../../../include/open-logic-section]{subfiles}

\begin{document}

\olfileid{mth}{ind}{idf}

\olsection{આગમનાત્મક વ્યાખ્યાઓ}

તર્કશાસ્ત્રમાં આપણે ઘણી વાર પદાર્થોના પ્રકારોને
\emph{આગમનાત્મક રીતે} વ્યાખ્યાયિત કરીએ છીએ. એટલે
જે પ્રકારને વ્યાખ્યાયિત કરવો હોય તેમાં પદાર્થ ક્યારે
ગણાય તેના નિયમો આપીએ છીએ; એ નિયમો તે પ્રકારના
અગાઉના પદાર્થોમાંથી નવા પદાર્થો કેવી રીતે બનાવવા
તે પણ સમજાવે છે. દાખલા તરીકે, ભાષાનાં પદો અને
!!{formula}s જેવી ચિહ્નોની ખાસ જાતની શ્રેણીઓ આપણે
ઘણી વાર આગમનથી વ્યાખ્યાયિત કરીએ છીએ. એક સરળ
ઉદાહરણ તરીકે, $\mathrm{a}$, $\mathrm{b}$, $\mathrm{c}$,
$\mathrm{d}$ અક્ષરો,~$\circ$ ચિહ્ન અને $[$ તથા $]$
કૌંસથી બનેલી શ્રેણીઓ વિચારો; જેમ કે
"$[[\mathrm{c} \circ \mathrm{d}][$",
"$[\mathrm{a}[]\circ]$",
"$\mathrm{a}$" અથવા "$[[\mathrm{a} \circ \mathrm{b}]\circ
\mathrm{d}]$". પહેલી બે શ્રેણીઓમાં કંઈક "ખોટું"
લાગે: પહેલીમાં કૌંસોનો જરાય મેળ બેસતો નથી;
અને તમને એમ પણ લાગે કે "$\circ$"એ પોતે અર્થપૂર્ણ
હોય એવી અભિવ્યક્તિઓને "જોડવી" જોઈએ. ત્રીજી
અને ચોથી શ્રેણી વધુ યોગ્ય લાગે છે: દરેક "$[$"
માટે બંધ કરતો "$]$" છે (જો કૌંસો હોય તો), અને
દરેક $\circ$ની બંને બાજુ કૌંસોની જોડીથી ઘેરાયેલી
"યોગ્ય" અભિવ્યક્તિઓ મળે છે.

"સુઘડ પદ" ક્યારે ગણાય તે આપણે ચોક્કસ રીતે કહેવા
માંગીએ છીએ. પહેલું તો, દરેક અક્ષર પોતે સુઘડ છે.
ફક્ત એક અક્ષર ન હોય એવું કંઈ પણ "$[t \circ s]$"
સ્વરૂપનું હોવું જોઈએ, જ્યાં $s$ અને $t$ પોતે
સુઘડ હોય. ઊલટું, $t$ અને $s$ સુઘડ હોય તો તેમની
વચ્ચે $\circ$ મૂકીને અને આસપાસ કૌંસોની જોડી
મૂકીને નવું સુઘડ પદ બનાવી શકીએ. આવી ક્રિયાઓથી
સુઘડ પદોના ગણને \emph{વ્યાખ્યાયિત} કરી શકીએ.
આ \emph{આગમનાત્મક
  વ્યાખ્યા} છે.

\begin{defn}[સુઘડ પદો]
  \emph{સુઘડ પદો}નો ગણ નીચે મુજબ આગમનાત્મક રીતે
  વ્યાખ્યાયિત થાય છે:
  \begin{enumerate}
  \item કોઈ પણ અક્ષર $\mathrm{a}$, $\mathrm{b}$, $\mathrm{c}$,
    $\mathrm{d}$ સુઘડ પદ છે.
  \item $s_1$ અને $s_2$ સુઘડ પદો હોય, તો
    $[s_1 \circ s_2]$ પણ સુઘડ પદ છે.
  \item આ સિવાય બીજું કંઈ સુઘડ પદ નથી.
  \end{enumerate}
\end{defn}

આ વ્યાખ્યા કહે છે કે કોઈ વસ્તુ સુઘડ પદ ત્યારે અને
ત્યારે જ છે જ્યારે શરતો (1) અને~(2) મુજબ તેને
સીમિત સંખ્યાનાં પગલાંમાં બનાવી શકાય. પહેલા પગલામાં
ફક્ત એક-એક અક્ષરથી બનેલાં બધાં સુઘડ પદો બનાવીએ:
\[
\mathrm{a}, \mathrm{b}, \mathrm{c}, \mathrm{d}
\]
બીજા પગલામાં બનાવેલા પદો પર (2) લાગુ કરીએ. બે
અક્ષરોની બધી જોડીઓમાંથી મળે:
\[
[\mathrm{a} \circ \mathrm{a}], [\mathrm{a} \circ \mathrm{b}],
[\mathrm{b} \circ \mathrm{a}], \dots, [\mathrm{d} \circ \mathrm{d}]
\]
ત્રીજા પગલામાં અત્યાર સુધી બનાવેલાં કોઈ પણ બે
સુઘડ પદો પર ફરી (2) લાગુ કરીએ. તેથી
$[\mathrm{a} \circ [\mathrm{a} \circ
    \mathrm{a}]]$ જેવા નવા સુઘડ પદો મળે: અહીં $t$
પહેલા પગલાનું $\mathrm{a}$ છે અને $s$ બીજા પગલાનું
$[\mathrm{a} \circ \mathrm{a}]$ છે. વળી
$[[\mathrm{b} \circ
    \mathrm{c}] \circ [\mathrm{d} \circ \mathrm{b}]]$
એ બીજા પગલાનાં બે પદો $[\mathrm{b} \circ \mathrm{c}]$
અને $[\mathrm{d}
  \circ \mathrm{b}]$માંથી બને છે. આ રીતે આગળ
વધીએ. ખંડ (3) ખાતરી કરે છે કે આ રીતે ન બનેલું
કશું સુઘડ પદોના ગણમાં ઘૂસી ન જાય.

ધ્યાન આપો કે જે દરેક ચિહ્ન-શ્રેણી આપણને "સુઘડ"
લાગે છે તે આ વ્યાખ્યા મુજબ પણ સુઘડ છે એવું
આપણે હજી સાબિત કર્યું નથી. પરંતુ જે કંઈ આ
રીતે બનાવી શકીએ તે ખરેખર "સુઘડ લાગે" છે
એ સ્પષ્ટ હોવું જોઈએ: કૌંસોનો મેળ બેસે છે અને
$\circ$ પોતે સુઘડ હોય એવા ભાગોને જોડે છે.

આગમનાત્મક વ્યાખ્યાઓનું મુખ્ય લક્ષણ એ છે કે જો
બધાં સુઘડ પદો વિશે કંઈ સાબિત કરવું હોય, તો વ્યાખ્યા
કયા કિસ્સા વિચારવાના છે તે કહે છે. દાખલા તરીકે,
$t$ સુઘડ પદ છે એમ કહીએ ત્યારે આગમનાત્મક
વ્યાખ્યા તેનું શક્ય સ્વરૂપ જણાવે છે: $t$ અક્ષર
હોઈ શકે અથવા કોઈ સુઘડ પદો $s_1$ અને~$s_2$
માટે $t$ એ $[s_1 \circ s_2]$ હોઈ શકે. ખંડ (3)ને
લીધે આ જ બે શક્યતાઓ છે.

આગમનાત્મક રીતે વ્યાખ્યાયિત ગણની બધી વસ્તુઓ
વિશેના દાવા સાબિત કરતી વખતે આગમનનું પ્રબળ સ્વરૂપ
ખાસ મહત્વનું બને છે. દાખલા તરીકે, ધારો કે
લંબાઈ~$n$ના દરેક સુઘડ પદમાં $[$ની સંખ્યા~$< n/2$
છે એવું સાબિત કરવું છે. તેને બધા~$n$ વિશેના
દાવા તરીકે જોઈ શકાય છે: દરેક $n$ માટે લંબાઈ~$n$ના
કોઈ પણ સુઘડ પદમાં $[$ની સંખ્યા $< n/2$ છે.

\begin{prop}
  કોઈ પણ $n$ માટે લંબાઈ~$n$ના સુઘડ પદમાં
  $[$ની સંખ્યા $< n/2$ છે.
\end{prop}

\begin{proof}
આ પરિણામ (પ્રબળ) આગમનથી સાબિત કરવા નીચેનો
શરતી દાવો સાચો છે એમ બતાવવું પડે:
\begin{quote}
  જો દરેક $l < k$ માટે લંબાઈ~$l$ના કોઈ પણ
  સુઘડ પદમાં $[$ની સંખ્યા $< l/2$ હોય,
  તો લંબાઈ~$k$ના કોઈ પણ સુઘડ પદમાં
  $[$ની સંખ્યા $< k/2$ હોય.
\end{quote}
આ શરતી દાવો બતાવવા તેનું પૂર્વપદ સાચું ધારો:
કોઈ પણ $l<k$ માટે લંબાઈ~$l$ના સુઘડ પદોમાં
$[$ની સંખ્યા $< l/2$ છે. આ ધારણાને
આગમન-પરિકલ્પના કહીએ છીએ. લંબાઈ~$k$ના
સુઘડ પદો માટે પણ એ જ વાત બતાવવી છે.

ધારો કે $t$ એ લંબાઈ~$k$નું સુઘડ પદ છે.
સુઘડ પદો આગમનાત્મક રીતે વ્યાખ્યાયિત હોવાથી
બે કિસ્સા છે: (1)~$t$ પોતે એક અક્ષર છે,
અથવા (2)~કોઈ સુઘડ પદો $s_1$ અને~$s_2$
માટે $t$ એ $[s_1 \circ s_2]$ છે.
\begin{enumerate}
\item $t$ એક અક્ષર છે. ત્યારે $k = 1$ અને
$t$માં $[$ની સંખ્યા~$0$ છે. $0 < 1/2$
હોવાથી દાવો સાચો છે.
\item કોઈ સુઘડ પદો $s_1$ અને~$s_2$ માટે
$t$ એ $[s_1 \circ s_2]$ છે. $l_1$ એ~$s_1$ની
અને $l_2$ એ~$s_2$ની લંબાઈ માનો. ત્યારે $t$ની
લંબાઈ~$k$ એ $l_1+l_2+3$ છે ($s_1$ અને
$s_2$ની લંબાઈ ઉપરાંત $[$, $\circ$, $]$
એવાં ત્રણ ચિહ્નો). $l_1+l_2+3$ હંમેશાં
$l_1$થી મોટું હોવાથી $l_1 < k$. એ જ રીતે
$l_2 < k$. તેથી આગમન-પરિકલ્પના $s_1$ અને~$s_2$
બંને પદો પર લાગુ પડે છે: $s_1$માં $[$ની
સંખ્યા~$m_1$ એ $< l_1/2$ છે અને~$s_2$માં
$[$ની સંખ્યા~$m_2$ એ $< l_2/2$ છે.

  $t$માં $[$ની સંખ્યા એ~$s_1$માંની $[$ની
  સંખ્યા, વત્તા~$s_2$માંની $[$ની સંખ્યા,
  વત્તા~$1$ છે; એટલે $m_1 + m_2 + 1$.
  $m_1 < l_1/2$ અને $m_2 < l_2/2$
  હોવાથી મળે:
  \[
  m_1 + m_2 + 1 < \frac{l_1}{2} + \frac{l_2}{2} + 1 = \frac{l_1+l_2+2}{2} < \frac{l_1+l_2+3}{2} = k/2.
  \]
\end{enumerate}
દરેક કિસ્સામાં આપણે (આગમન-પરિકલ્પનાના આધારે)
બતાવ્યું કે $t$માં $[$ની સંખ્યા $< k/2$ છે.
પ્રબળ આગમનથી પ્રતિજ્ઞાપન ફલિત થાય છે.
\end{proof}

\begin{prob}
  અતિસુઘડ પદોનો ગણ આ રીતે વ્યાખ્યાયિત કરો:
  \begin{enumerate}
  \item કોઈ પણ અક્ષર $\mathrm{a}$, $\mathrm{b}$, $\mathrm{c}$,
    $\mathrm{d}$ અતિસુઘડ પદ છે.
  \item $s$ અતિસુઘડ પદ હોય તો $[s]$ પણ છે.
  \item $s_1$ અને $s_2$ અતિસુઘડ પદો હોય, તો
    $[s_1 \circ s_2]$ પણ છે.
  \item આ સિવાય બીજું કંઈ અતિસુઘડ પદ નથી.
  \end{enumerate}
  લંબાઈ~$n$ના અતિસુઘડ પદ~$t$માં $[$ની
  સંખ્યા $\le n/2 +1$ છે એમ બતાવો.
\end{prob}

\end{document}
